\chapter{Einleitung} \label{ch:Einleitung}

Die Ausarbeitung der Einleitung ist anhand des Skriptes von Jens Wagner entstanden. Alle Abbildungen, sind entweder aus dem Skript oder selbst erzeugt, wenn nicht anders vermerkt. Rechtschreibprüfungen werden sowohl durch K.I.-Modelle, sowie durch den Rechtschreibprüfer des Dudens durchgeführt.  Dies gilt nur für die Einleitung und die Durchführung. Die Auswertung wird völlig ohne K.I. Assistenz geschrieben. \cite{skriptPAP21, Duden}


\section{Aufgabe und Motivation} \label{sec:AuM}

Gekoppelte Oszillatoren spielen in vielen Bereichen der Physik eine zentrale Rolle. Sie treten sowohl in mechanischen Systemen wie Pendeln, als auch in Molekülen, elektrischen Schwingkreisen oder Festkörpern auf. Ein besonders anschauliches Beispiel ist das System zweier mechanisch gekoppelter Pendel, an dem sich grundlegende Eigenschaften wie Normalschwingungen, Energieübertragung und Schwebung sehr gut beobachten lassen.

Ziel dieses Versuchs ist es, die Dynamik zweier über eine Feder gekoppelter Pendel experimentell zu untersuchen. Dazu werden die Eigenfrequenzen der symmetrischen und antisymmetrischen Schwingung bestimmt, die Schwebung bei gemischter Anregung analysiert und der Kopplungsgrad quantitativ ermittelt. Ergänzend wird das Verhalten eines analogen elektrischen Systems betrachtet, um die Übertragbarkeit der Konzepte auf andere physikalische Bereiche zu verdeutlichen.


\section{Physikalische Grundlage} \label{sec:PhyBasics}

\subsection{Das einfache Pendel}

Für kleine Auslenkungen lässt sich die Bewegung eines einzelnen physikalischen Pendels durch die linearisierte Differentialgleichung
\eq{einzelpendel}{
    J\ddot{\varphi} = -D\varphi
}
beschreiben. Hierbei ist \(J\) das Trägheitsmoment des Pendels und \(D = mgL\) das Direktionsmoment, wobei \(m\) die Masse, \(g\) die Erdbeschleunigung und \(L\) die Pendellänge ist. Die Lösung dieser Gleichung ist eine harmonische Schwingung mit der Eigenkreisfrequenz
\eq{omega_einzel}{
    \omega = \sqrt{\frac{D}{J}}.
}

\subsection{Das gekoppelte Pendel}

Werden zwei identische Pendel durch eine Feder gekoppelt, so wirkt zusätzlich ein kopplungsbedingtes Drehmoment. Das Direktionsmoment der Feder lautet
\eq{direktionsmoment_feder}{
    D' = D_\mathrm{F} \, l^2,
}
wobei \(D_\mathrm{F}\) die Federkonstante und \(l\) der Abstand zwischen Federaufhängung und Pendelachse ist. Die Bewegungsgleichungen der beiden Pendel lauten dann
\eq{gekoppelte_DGL}{
\begin{aligned}
    J\ddot{\varphi}_1 &= -D\varphi_1 + D'(\varphi_2 - \varphi_1), \\
    J\ddot{\varphi}_2 &= -D\varphi_2 + D'(\varphi_1 - \varphi_2).
\end{aligned}
}

Durch die Substitution \(u = \varphi_1 + \varphi_2\) und \(v = \varphi_1 - \varphi_2\) entkoppelt sich das Gleichungssystem in zwei unabhängige Differentialgleichungen:
\eq{entkoppelt}{
\begin{aligned}
    J\ddot{u} + D u &= 0, \\
    J\ddot{v} + (D + 2D') v &= 0.
\end{aligned}
}

\subsection{Symmetrische Oszillation}
\wrap{R}{memes/balanced}{Viel geiler als Schwebung...}
Bei der symmetrischen Schwingung werden beide Pendel gleich weit und gleichsinnig ausgelenkt. Die Kopplungsfeder wird dabei nicht gedehnt, weshalb die Schwingung mit der Eigenfrequenz des freien Pendels erfolgt:
\eq{omega_sym}{
    \omega_1 = \sqrt{\frac{D}{J}}.
}
Die Pendel schwingen phasengleich und behalten stets denselben Abstand zueinander.


\subsection{Anti-Symmetrische Oszillation}

Werden die Pendel gegensinnig ausgelenkt, so wird die Kopplungsfeder maximal beansprucht. Dadurch erhöht sich das rücktreibende Moment, und die resultierende Eigenfrequenz ist größer:
\eq{omega_asym}{
    \omega_2 = \sqrt{\frac{D + 2D'}{J}}.
}
Die Pendel schwingen dabei exakt gegenphasig.

\subsection{Schwebung}

Werden die Pendel nicht in einer Normalschwingung angeregt, sondern beispielsweise nur eines ausgelenkt, so überlagern sich beide Eigenmoden. Dies führt zu einer periodischen Energieübertragung zwischen den Pendeln, der sogenannten Schwebung. Die Einzelfrequenz und die Schwebungsfrequenz ergeben sich zu
\eq{schwebung}{
    \omega_\mathrm{I} = \frac{\omega_1 + \omega_2}{2}, \qquad
    \omega_\mathrm{II} = \frac{\omega_2 - \omega_1}{2}.
}
Die Schwebungsfrequenz beschreibt die Rate, mit der die Energie zwischen den Pendeln hin- und herwandert.

\subsection{Die Kopplungsstärke}

Zur Quantifizierung der Kopplung wird der Kopplungsgrad definiert als
\eq{kopplungsgrad}{
    K = \frac{D'}{D + D'}.
}
Über die Eigenfrequenzen lässt sich dieser experimentell bestimmen:
\eq{kopplungsgrad_frequenzen}{
    K = \frac{\omega_2^2 - \omega_1^2}{\omega_2^2 + \omega_1^2}.
}
Für schwache Kopplungen (\(D' \ll D\)) gilt näherungsweise
\eq{kopplungsgrad_naeherung}{
    K \approx \frac{D'}{D} \propto l^2,
}
wodurch sich die Kopplungsgrade verschiedener Federaufhängungen direkt über die Abstände vergleichen lassen.


\subsection{Der Hall-Effekt}
Der im Versuchsaufbau verwendete Hallsensor dient zur präzisen Messung der Winkelauslenkung der Pendel. Er befindet sich an der Pendelachse und bewegt sich relativ zu einem feststehenden Magnetfeld. Durch den Hall-Effekt entsteht dabei eine Spannung, deren Betrag proportional zum Sinus des Auslenkwinkels ist. Diese Spannung wird verstärkt, digitalisiert und anschließend von der Messsoftware ausgewertet. Auf diese Weise lässt sich die zeitabhängige Pendelbewegung berührungslos, kontinuierlich und mit hoher Genauigkeit erfassen, was die Grundlage für die spätere Frequenzanalyse bildet.

% \subsubsection*{Skizze des Versuchsaufbaus} \label{sec:Skizze}
